\section{Proofs for finite-prize expected utility}
\label{sec:annex-p17-finite-expected-utility}

\begin{annexitem}{annex:p17.lottery-reduction}{thm:p17-lottery-reduction}{The probability law of a finite mixture}
\begin{proof}
For every prize \(z\), each summand \(\alpha_kp_k(z)\) is nonnegative.
Finite summation gives
\[
 \sum_{z\in Z}\sum_{k=1}^K\alpha_kp_k(z)
   =\sum_{k=1}^K\alpha_k\sum_{z\in Z}p_k(z)
   =\sum_{k=1}^K\alpha_k=1.
\]
No coordinate exceeds one because all coordinates are nonnegative and
their sum is one. Thus the mixture is a lottery. For any fixed \(w\in Z\),
the coordinate of \(\sum_zp(z)\delta_z\) at \(w\) is \(p(w)\), since
all terms except the one at \(z=w\) vanish. This proves the decomposition.

For the nested mixture, its coordinate at \(z\) is
\[
 \sum_k\alpha_k\left(\sum_\ell\beta_{k\ell}p_{k\ell}(z)\right)
       =\sum_k\sum_\ell\alpha_k\beta_{k\ell}p_{k\ell}(z).
\]
Distributivity for finite sums proves the identity at each coordinate.
The flattened weights are nonnegative and sum to
\(\sum_k\alpha_k\sum_\ell\beta_{k\ell}=1\); no infinite sum or
limiting interchange is involved.

On the finite set \(\{1,\ldots,K\}\times Z\), assign mass
\(\alpha_kp_k(z)\) to \((k,z)\). Its total mass is one by the first
calculation. Summing over \(z\) gives branch marginal \(\alpha_k\),
and summing over \(k\) gives the stated final prize law. If
\(\alpha_k>0\), the conditional prize mass in that branch is
\(\alpha_kp_k(z)/\alpha_k=p_k(z)\). If \(\alpha_k=0\), every joint
mass in that branch is zero and no conditional ratio is taken. Any
specified \(p_k\) on it leaves the final law unchanged. This proves
the two-stage claim with its null-branch convention.
\end{proof}
\end{annexitem}

\begin{annexitem}{annex:p17.mixture-comparison}{thm:p17-mixture-comparison}{Component replacement and strict mixture order}
\begin{proof}
We first record elementary consequences of completeness and transitivity
that apply even though the lottery domain is infinite. Indifference is
reflexive and symmetric by definition, and is transitive by applying weak
transitivity in both directions. If \(p\sim p'\) and \(q\sim q'\),
then \(p\succeq q\) implies \(p'\succeq p\succeq q\succeq q'\);
the reverse implication is the same chain in the other direction.
Thus replacing indifferent objects preserves comparisons. Also,
\(p\succeq q\succ r\) implies \(p\succ r\): weak transitivity gives
\(p\succeq r\), while \(r\succeq p\) would imply
\(r\succeq q\), contradicting \(q\succ r\). The argument for
\(p\succ q\succeq r\) is analogous. A finite chain of weak comparisons
with a strict step is consequently strict between its endpoints.

For \(0<\alpha<1\), independence gives equivalence of the forward
weak comparisons of \(p,q\) and their two mixtures. Applying the same
axiom with \(p,q\) exchanged gives equivalence of the reverse weak
comparisons. Combining these two facts proves preservation of strict
preference and of indifference.

Consider the replacement of the \(j\)-th component of a finite mixture,
with weight \(\alpha_j\). If \(\alpha_j=0\), the vector is unchanged.
If \(\alpha_j=1\), all other weights vanish and the comparison is
directly the comparison of the replaced lotteries. For
\(0<\alpha_j<1\), form the common residual lottery
\[
             r=\sum_{k\ne j}\frac{\alpha_k}{1-\alpha_j}p_k.
\]
The residual weights are nonnegative and sum to one. The old and new
mixtures are \(\alpha_jp_j+(1-\alpha_j)r\) and
\(\alpha_jq_j+(1-\alpha_j)r\). The axiom and its strict/indifference
consequences establish the required comparison. The separate weight-one
case ensures that this proof never divides by zero or needs a lottery
formed from an empty residual family. Applying this argument successively
to the finitely many components, and using the relation facts above,
proves the finite replacement assertion.

Suppose now that \(a\succ b\). For \(0<t<1\), independence applied
to \(a,b\) with common component \(b\) and weight \(t\) gives
\(B_t\succ b\). Applying it with common component \(a\) and weight
\(1-t\) gives
\[
 (1-t)a+ta=a\succ(1-t)b+ta=B_t.
\]
Thus \(a\succ B_t\succ b\). These inequalities also give the
comparison of \(B_s,B_t\) whenever \(s=1\) or \(t=0\), with
the case \(s=1,t=0\) being the assumed \(a\succ b\).

For \(0<t<s<1\), let \(\lambda=t/s\), so \(0<\lambda<1\).
We have already established \(B_s\succ b\). Independence, with
common component \(B_s\) and weight \(1-\lambda\), yields
\[
 B_s=(1-\lambda)B_s+\lambda B_s
       \succ(1-\lambda)b+\lambda B_s=B_t.
\]
The last equality follows at each coordinate from \(\lambda s=t\).
This covers every strict ordering of two parameters in the closed interval.

Finally, if \(a\sim b\), indifference preservation with common
component \(b\) shows \(B_t\sim b\) for \(0<t<1\). The endpoints
are \(b,a\), already indifferent; transitivity supplies indifference
to both throughout. This completes all cases.
\end{proof}
\end{annexitem}

\begin{annexitem}{annex:p17.lottery-calibration}{thm:p17-lottery-calibration}{Existence and uniqueness of the calibration probability}
\begin{proof}
Write \(B_t=ta+(1-t)b\), and define
\[
                 A=\{t\in[0,1]:p\succeq B_t\}.
\]
The set contains zero because \(p\succeq b\), and is bounded above
by one. For \(s,t\in[0,1]\), direct computation gives
\[
                      d(B_s,B_t)=|s-t|d(a,b).
\]
Here \(d(a,b)>0\), since \(a\succ b\) excludes equality of the
two vectors. Thus the segment map is continuous.

The set \(A\) is closed relative to \([0,1]\). Indeed, if
\(t\notin A\), completeness gives \(B_t\succ p\). Continuity
at that strict comparison, keeping its second argument \(p\) fixed,
implies \(B_s\succ p\) for every \(s\) sufficiently close to \(t\).
Those \(s\) also lie outside \(A\), so its complement is relatively
open. This argument uses exactly the neighborhood continuity in the body.

Let \(t_* =\sup A\), which exists by completeness of \(\mathbb R\).
Then \(t_*\in A\). To see this directly, when \(t_*=0\) use
\(0\in A\); otherwise the supremum property supplies points of \(A\)
arbitrarily close to \(t_*\) from below, and closedness includes their
limit. Hence \(p\succeq B_{t_*}\).

If \(B_{t_*}\not\succeq p\), this last comparison is strict. If
\(t_*=1\), it contradicts the hypothesis \(a\succeq p\), since
\(B_1=a\). If \(t_*<1\), continuity of \(p\succ B_{t_*}\) and
the segment map give some \(s\) with \(t_*<s\leq1\) and
\(p\succ B_s\). Such an \(s\) belongs to \(A\), contradicting
that \(t_*\) is an upper bound. Therefore \(B_{t_*}\succeq p\)
and \(p\sim B_{t_*}\).

If two parameters \(s>t\) both gave indifference to \(p\),
transitivity of indifference would give \(B_s\sim B_t\). The strict
segment order in \cref{thm:p17-mixture-comparison} instead gives
\(B_s\succ B_t\). Thus the parameter is unique. Since the segment
endpoints are \(b,a\), uniqueness identifies the parameter with zero
or one exactly in the corresponding endpoint-indifference cases. If
\(a\succ p\succ b\), neither endpoint is possible, so the parameter
is interior. Conversely an interior parameter is strictly between the
endpoints by the segment lemma and indifference replacement. This also
proves the stated strict interpolation property.
\end{proof}
\end{annexitem}

\begin{annexitem}{annex:p17.finite-vnm}{thm:p17-finite-vnm}{Construction and converse of the finite-prize representation}
\begin{proof}
Assume first the four stated preference conditions. Restrict the weak
preference to the finitely many certain lotteries. This restriction
has a best and a worst member: enumerate the finite nonempty prize set,
retain a weakly better member after each comparison to find a best one,
and retain a weakly worse member to find a worst one. Completeness
makes each selection possible; transitivity ensures that the retained
member compares in the required direction with all earlier members.
Denote the resulting certain lotteries by \(M,m\). For each prize,
\(M\succeq\delta_z\succeq m\).

Every lottery \(p\) lies weakly between \(M\) and \(m\).
In the decomposition \(p=\sum_zp(z)\delta_z\), replace the certain
components successively by \(M\). The component-replacement lemma
gives \(M\succeq p\), because the resulting mixture is
\(\sum_zp(z)M=M\). Replacing them by \(m\) instead gives
\(p\succeq m\). This includes lotteries with zero coordinates.

If all certain prizes are indifferent, replace every component by
\(m\) using indifferent replacements. Then every \(p\sim m\),
so every two lotteries are indifferent. The constant assignment
\(u(z)=0\) gives exactly that relation by expectation. A singleton
prize set has only one certain lottery and only that lottery in its
simplex, so is included in this argument.

Otherwise \(M\succ m\). For if \(M\sim m\), the inequalities
\(M\succeq\delta_z\succeq m\) would make every certain prize
indifferent by transitivity. Apply \cref{thm:p17-lottery-calibration}
to each \(\delta_z\), and let its unique calibration parameter be
\(u(z)\). This defines a function on all of \(Z\), with values in
\([0,1]\), value one at the chosen best prize and zero at the chosen
worst prize. Ties with either endpoint are included by the calibration
lemma.

For each lottery \(p\), replace the \(z\)-th component in its certain-prize
decomposition by its indifferent binary lottery
\(u(z)M+(1-u(z))m\). Finite indifferent replacement gives
\[
\begin{aligned}
 p&\sim\sum_zp(z)\bigl(u(z)M+(1-u(z))m\bigr)\\
  &=\left(\sum_zp(z)u(z)\right)M
          +\left(1-\sum_zp(z)u(z)\right)m.
\end{aligned}
\]
The equality is the finite mixture-reduction identity. The coefficient
\(U_u(p)=\sum_zp(z)u(z)\) is in \([0,1]\) because it is an average
of numbers in that interval. The analogous identity holds for \(q\).
Indifference replacement and the strict segment order therefore give
\[
 p\succeq q
  \quad\Longleftrightarrow\quad
  U_u(p)M+(1-U_u(p))m\succeq U_u(q)M+(1-U_u(q))m
  \quad\Longleftrightarrow\quad U_u(p)\geq U_u(q).
\]
For equal coefficients the mixtures are identical, and for unequal
coefficients their strict order was proved, so both directions and ties
are covered. Uniqueness under the stated normalization follows either
from uniqueness of each calibration probability or from the affine
uniqueness theorem proved next, without using that later result here.

Conversely, suppose a real function \(u\) represents the relation by
the finite sum \(U_u\). The usual order of real numbers makes this
relation complete and transitive. Distributivity gives, for every
\(\alpha\in[0,1]\),
\[
 U_u(\alpha p+(1-\alpha)r)=\alpha U_u(p)+(1-\alpha)U_u(r).
\]
For \(0<\alpha<1\), subtracting the two common-component expressions
leaves \(\alpha(U_u(p)-U_u(q))\), whose sign is the sign of the
unmixed difference. This proves the independence equivalence.

For continuity put \(C=\sum_z|u(z)|\). For any lotteries \(p,p'\),
\[
 |U_u(p')-U_u(p)|
   \leq\sum_z|u(z)|\,|p'(z)-p(z)|
   \leq C d(p',p),
\]
because each coordinate difference has absolute value at most the
Euclidean distance. If \(p\succ q\), set
\(g=U_u(p)-U_u(q)>0\). This forces \(C>0\). Choose
\(\varepsilon=g/(4C)\). When both lottery arguments are within
\(\varepsilon\) of \(p,q\), their score difference is greater than
\(g-2C\varepsilon=g/2>0\). This is the required strict-neighborhood
condition. If there is no strict pair, continuity is vacuous; no division
by zero is performed. This proves the converse and the theorem.
\end{proof}
\end{annexitem}

\begin{annexitem}{annex:p17.vnm-affine-uniqueness}{thm:p17-vnm-affine-uniqueness}{The consequence scale determined by lottery comparisons}
\begin{proof}
First suppose \(u\) is nonconstant. Choose prizes \(z^+,z^-\) at
which it attains its largest and smallest values, respectively. Write
\(u_+=u(z^+)\), \(u_-=u(z^-)\), and similarly \(v_+,v_-\).
We have \(u_+>u_-\), hence \(\delta_{z^+}\succ\delta_{z^-}\).
Representation by \(v\) gives \(v_+>v_-\). Define
\[
                  a=\frac{v_+-v_-}{u_+-u_-}>0,
                  \qquad b=v_- -a u_-.
\]
For any prize \(z\), let
\(t=(u(z)-u_-)/(u_+-u_-)\), which is in \([0,1]\). The lottery
\(t\delta_{z^+}+(1-t)\delta_{z^-}\) has \(u\)-expectation
\(t u_++(1-t)u_-=u(z)\), so it is indifferent to \(\delta_z\).
The \(v\)-representation must give equal values on this indifference:
\[
\begin{aligned}
 v(z)&=t v_++(1-t)v_-\\
     &=t(a u_++b)+(1-t)(a u_-+b)
       =a u(z)+b.
\end{aligned}
\]
The distinct endpoint \(u\)-values determine these two constants
uniquely. This argument uses equality of the represented preference on
all lotteries, rather than merely equality of its ranking on prizes.

If \(u\) is constant with value \(c\), every lottery has expected
value \(c\), so all are indifferent. In particular all certain lotteries
are indifferent under the \(v\)-representation, forcing \(v(z)=d\)
for a single constant \(d\). For any \(a>0\), choosing \(b=d-ac\)
then gives \(v=a u+b\) on the entire nonempty prize set.

Finally, for any \(v=a u+b\) with \(a>0\), normalization of a
lottery gives
\[
 U_v(p)=\sum_zp(z)(a u(z)+b)
       =aU_u(p)+b\sum_zp(z)=aU_u(p)+b.
\]
Its positive slope preserves all weak and strict comparisons and all
ties. This proves the converse including constant utilities and
singleton prize sets.
\end{proof}
\end{annexitem}

\begin{annexitem}{annex:p17.vnm-counterpreferences}{thm:p17-vnm-counterpreferences}{Explicit failures of the two lottery axioms}
\begin{proof}
The largest-probability criterion is represented by the real score
\(V(p)=\max\{p(a),p(b)\}\), so completeness and transitivity follow
from the real order. For any \(p,q\),
\[
                   |V(p)-V(q)|\leq\max_{z\in\{a,b\}}|p(z)-q(z)|
                                        \leq d(p,q).
\]
For the first inequality, choose a coordinate attaining \(V(p)\)
to obtain \(V(p)-V(q)\leq\max_z|p(z)-q(z)|\); exchange \(p,q\)
for the other direction. A positive gap \(V(p)-V(q)=g\) persists
when each argument moves by less than \(g/4\), proving continuity.
But both certain lotteries have value one and are indifferent. Mixing
each with the common \(\delta_a\) at weight \(1/2\) leaves the
first certain and makes the second uniform, with values one and one-half.
Indifference preservation, which follows from independence, fails.

For the lexicographic relation, distinct vectors have a first differing
coordinate, so they are compared, while equal vectors are indifferent.
For transitivity of strict comparisons \(p\succ q\succ r\), let
\(i,j\) be the first differing coordinates of the respective pairs.
Before \(\min(i,j)\) all three vectors coincide. If \(i<j\),
then \(p_i>q_i=r_i\); if \(j<i\), then \(p_j=q_j>r_j\);
if \(i=j\), then \(p_i>q_i>r_i\). Each case gives \(p\succ r\).
Comparisons involving equality preserve the other comparison directly,
so the weak relation is transitive as well as complete.

In mixtures with a common \(r\) and positive weight \(\alpha\),
the difference of the two vectors at coordinate \(z\) is
\(\alpha(p(z)-q(z))\). The first unequal coordinate is unchanged,
and its sign is unchanged because \(\alpha>0\). Equal vectors remain
equal. Thus both directions of the independence axiom hold, including
all weak comparisons.

In the body's three-coordinate example, \(p,q\) tie in their first
coordinate and \(p\)'s second coordinate is larger, so \(p\succ q\).
For \(0<\varepsilon<1/2\), the vector \(q_\varepsilon\) is a
lottery and its first coordinate is larger than that of \(p\), so
\(q_\varepsilon\succ p\). Its distance from \(q\) is
\(\sqrt2\,\varepsilon\). Given any positive neighborhood radius,
choose such an \(\varepsilon\) small enough to put
\(q_\varepsilon\) within that radius of \(q\), while leaving
\(p\) fixed. The pair contradicts the strict-neighborhood requirement
for \(p\succ q\). Continuity therefore fails.

An expected-utility representation would satisfy both axioms by the
already proved converse in \cref{thm:p17-finite-vnm}. Each exhibited
failure precludes that representation on its stated domain. This is a
mathematical exclusion and makes no claim about which preference a
particular person has.
\end{proof}
\end{annexitem}
